\subsubsection*{1.4.1 Business Narrative}
\addcontentsline{toc}{subsubsection}{1.4.1 Business Narrative}
\pdfbookmark[2]{1.4.1 Business Narrative}{p7h7}
Author: 722I0004 \ensuremath{\cdot} AI-assisted (Codex)

SaigonBite's hubs are delivery-only kitchens. Each kitchen prepares food under several online-only brands, and shippers carry orders to homes, offices or hotels. In the case's present-day process, customers contact SaigonBite by phone, Zalo or third-party delivery apps; staff rewrite orders onto paper tickets, while shippers are arranged by calls or group chat. Kitchen workload is hard to see outside each hub, menus can lag behind ingredient availability, cash handoffs can be mistaken, and discounts and refunds are handled manually. Customers and managers therefore lack one view of orders, availability, delivery performance and sales. The proposed platform is meant to coordinate those handoffs and make hub capacity, order progress, payment and reporting visible.

Customers are local residents or foreign tourists. A customer creates an account, keeps delivery addresses, browses hubs, virtual brands and hub-specific menu items, then submits one order for one cooking hub. Local registration requires a Vietnamese mobile number; tourist registration uses passport number, nationality and email, with phone optional and a hotel address allowed. Passport information is limited to customer and administrator views. The customer chooses an address, at most one optional voucher and a payment option. Before confirmation the platform shows a full quote: item subtotal, voucher discount, distance-based delivery fee, any quiet-hub fee reduction and total. This makes consent depend on the price the customer actually sees.

At checkout, workload and delivery range shape the hub choice. The plan counts PLACED, CONFIRMED and PREPARING orders as kitchen load. At capacity a hub refuses a new order; the system can suggest up to three active, non-busy alternatives in range, but their menus may differ and require a new quote. A quieter hub may apply a fee discount. The case leaves the figures open, so the current values---25\% of capacity, a 5,000 VND discount, an 8 km radius, distance estimated to a district centroid and 3,000 VND per started kilometre on top of the hub's base fee---are team defaults, not client answers. The quote labels the centroid distance as approximate.

The plan offers cash on delivery (COD) to all customers. Online choices are e-wallet or domestic bank QR for local customers and international card for tourists. In this release, those online outcomes use a team-built simulator; no real card data or provider processing is represented. A declined simulation creates no order. An approved online checkout is recorded as paid within the simulated boundary, while COD remains pending until the shipper records cash collection. The customer can open a printable receipt with the order and address snapshots, item lines, charges, total, payment type and state, and any refund owed. It is an order view, not a separate persisted transaction.

After an accepted checkout, the selected hub receives a PLACED order. Kitchen staff manage their own hub's queue: they may confirm or reject the order with a reason, then start preparation and mark it ready. They also update item availability. If a dish is unavailable while the order is still PLACED, staff may offer a different listed, available item at the same or lower unit price. The customer sees the proposed change and updated total before choosing. Acceptance updates the order; rejection or declining an offer closes it. Customers can track order progress and cancel while it remains PLACED or CONFIRMED; the plan blocks cancellation after preparation begins.

After the kitchen marks an order READY, a shipper takes it, records pickup, then records successful delivery or a failure. For COD, the shipper records cash collection equal to the payable order total. Rejection, cancellation, a declined replacement or delivery failure voids COD. After those terminal outcomes, online payment retains its original capture, enters REFUND\_PENDING and records the full captured amount as owed. If a customer accepts a cheaper replacement, the payment remains PAID while the price difference is recorded as refund owed. The system does not claim that money has been paid back. Complaint decisions, refund payout and rating remain deferred; L3 was accepted at the Part 1 review (no change requested).

Administrators maintain hubs and capacity, brands, menu items and hub listings, vouchers and user accounts. Closed records are archived so operational lists stay clear while authorised history remains available. Managers read sales, hub workload and delivery performance reports. The plan defines sales by successful delivery finish time, workload as a current snapshot, and delivery intervals and failure rate from lifecycle timestamps. These are operational reports, not accounting. Part 2 adds minimal supporting paths for registration, tracking and receipts, cancellation, preparation, availability, vouchers, account maintenance and order archive under the team's CRUD decision D-4. The defaults in the log remain open to the unanswered client questions and any future course clarification; L6 remains open. They are not client-approved business policy.

Source basis: RAD Sem01\_2627 Case Study, pp. 1--2; RAD Final Project Requirements, \S{}\S{}2.1--2.2 and 4; locked Part 2 plan, \S{}\S{}1--4; docs/part2/rules.md, BR-01--BR-26; docs/part2/decision\_log.md, D-2 and D-4.
